\par\Needspace{18\baselineskip}\bigskip
\pdfbookmark[2]{First section. Elliptic motion.}{pallas-II-section1}
\begin{center}
{\large\bfseries First section.\par}
\medskip\textit{Elliptic motion.}
\end{center}
\begin{flushright}
\begin{minipage}{.94\textwidth}
\small\raggedright
. . . . . Ye other wandring fires that move\\
In mystic dance, not without song, resound\\
His praise, who out of darkness call’d up light.\\
\hspace*{2em}\textsc{Milton}, Paradise lost.
\end{minipage}
\end{flushright}
\pdfbookmark[3]{Exposition, 1}{pallas-II-article1}
\begin{center}1.\end{center}
\noindent\hspace*{\parindent}Let \(x,y,z\) be the coordinates of a planet with respect to three mutually perpendicular planes passing through the center of the Sun. Let \(t\) be the time, \(m\) the sum of the masses of the Sun and the planet, and \(r\) the distance of the \SourceLeaf{447}planet from the Sun. Disregarding perturbations, we then have
\[
\begin{array}{rl}
1)\quad&0=\dfrac{\mathrm{dd}x}{\mathrm{d}t^2}+\dfrac{mx}{r^3}\\[5pt]
2)\quad&0=\dfrac{\mathrm{dd}y}{\mathrm{d}t^2}+\dfrac{my}{r^3}\\[5pt]
3)\quad&0=\dfrac{\mathrm{dd}z}{\mathrm{d}t^2}+\dfrac{mz}{r^3}\mathpunct{\cdot}
\end{array}
\]
\par It follows that if we set
\[
\begin{array}{rl}
4)\quad&\dfrac{y\mathrm{d}z-z\mathrm{d}y}{\mathrm{d}t}=A\\[5pt]
5)\quad&\dfrac{z\mathrm{d}x-x\mathrm{d}z}{\mathrm{d}t}=B\\[5pt]
6)\quad&\dfrac{x\mathrm{d}y-y\mathrm{d}x}{\mathrm{d}t}=C,
\end{array}
\]
\(A,B,C\) will be constant quantities and will therefore depend on the initial state. From these equations we obtain
\[
\begin{array}{rl}7)\quad&Ax+By+Cz=0.\end{array}
\]
\par The orbit therefore lies in a plane passing through the center of the Sun. We also have
\[
\begin{gathered}
rr=xx+yy+zz\\
r\mathrm{d}r=x\mathrm{d}x+y\mathrm{d}y+z\mathrm{d}z\\
\begin{aligned}
r\mathrm{dd}r+\mathrm{d}r^2&=x\mathrm{dd}x+y\mathrm{dd}y+z\mathrm{dd}z+\mathrm{d}x^2+\mathrm{d}y^2+\mathrm{d}z^2\\
&=\mathrm{d}x^2+\mathrm{d}y^2+\mathrm{d}z^2-\frac{m\mathrm{d}t^2}{r}\mathpunct{\cdot}
\end{aligned}
\end{gathered}
\]
\par Now equations 4, 5, 6 give
\[
\begin{aligned}
(AA+BB+CC)\mathrm{d}t^2
&=(xx+yy+zz)(\mathrm{d}x^2+\mathrm{d}y^2+\mathrm{d}z^2)\\
&\qquad{}-(x\mathrm{d}y+y\mathrm{d}y+z\mathrm{d}z)^2\\
&=rr(\mathrm{d}x^2+\mathrm{d}y^2+\mathrm{d}z^2)-rr\mathrm{d}r^2\\
&=r^3\mathrm{dd}r+mr\mathrm{d}t^2.
\end{aligned}
\]
\par Thus if we set
\[
\frac{AA+BB+CC}{m}=p,\qquad r-p=u,
\]
we obtain \(0=\frac{\mathrm{dd}u}{\mathrm{d}t^2}+\frac{u}{r^3}\mathpunct{\raisebox{.45ex}{,}}\) an equation analogous to the first three. We therefore see that if we set
\[
\begin{array}{rl}
8)\quad&\dfrac{x\mathrm{d}u-u\mathrm{d}x}{\mathrm{d}t}=F\\[5pt]
9)\quad&\dfrac{y\mathrm{d}u-u\mathrm{d}y}{\mathrm{d}t}=G\\[5pt]
10)\quad&\dfrac{z\mathrm{d}u-u\mathrm{d}z}{\mathrm{d}t}=H,
\end{array}
\]
\SourceLeaf{448}\(F,G,H\) will also be constant quantities, depending on the initial state. Combining these formulas with 4, 5, 6 gives
\[
\begin{array}{rl}
11)\quad&Gz-Hy=A(r-p)\\
12)\quad&Hx-Fz=B(r-p)\\
13)\quad&Fy-Gx=C(r-p).
\end{array}
\]
\par From these equations we readily conclude that the projections of the orbit and the orbit itself are curves of the second order, that the orbit has semiparameter \(p\), and that the center of the Sun occupies one of its foci. We note also that they immediately give an equation of condition relating the six constants \(A,B,C,F,G,H\)
\[
\begin{array}{rl}14)\quad&AF+BG+CH=0.\end{array}
\]
\par Similarly, from equations 8, 9, 10 we obtain
\[
\begin{aligned}
&(FF+GG+HH)\mathrm{d}t^2\\
&=(xx+yy+zz)\mathrm{d}u^2-2(x\mathrm{d}x+y\mathrm{d}y+z\mathrm{d}z)u\mathrm{d}u\\
&\qquad{}+(\mathrm{d}x^2+\mathrm{d}y^2+\mathrm{d}z^2)uu\\
&=rr\mathrm{d}r^2-2r\mathrm{d}r\mathpunct{.}(r-p)\mathrm{d}r+(\mathrm{d}x^2+\mathrm{d}y^2+\mathrm{d}z^2)(r-p)^2\\
&=(AA+BB+CC)\left(1-\frac{2p}{r}\right)\mathrm{d}t^2+(\mathrm{d}x^2+\mathrm{d}y^2+\mathrm{d}z^2)pp.
\end{aligned}
\]
\par Therefore if we set
\[
\begin{array}{rl}15)\quad&\dfrac{FF+GG+HH}{AA+BB+CC}=1-\dfrac{p}{a}\mathpunct{\raisebox{.45ex}{,}}\end{array}
\]
we obtain
\[
\begin{array}{rl}16)\quad&0=\dfrac{\mathrm{d}x^2+\mathrm{d}y^2+\mathrm{d}z^2}{\mathrm{d}t^2}-\dfrac{2m}{r}+\dfrac{m}{a}\mathpunct{\cdot}\end{array}
\]
\par We could have obtained this same equation by multiplying equations 1, 2, 3 by \(2\mathrm{d}x,\ 2\mathrm{d}y,\ 2\mathrm{d}z\) and integrating the sum. The preceding derivation has the advantage of showing the relation between the constant that this integration would have introduced and those already found.
\par
